%! Author = alexpayne
%! Date = 9/4/24

% Preamble
\documentclass[12pt, english, course]{lecture}

% Packages
\usepackage{amsmath}

\title{Mathematical Foundations for Clustering, Similarity, and Network Analysis}

% Document
\begin{document}
Two kinds of mathematical frameworks are discussed in this chapter: 
clustering algorithms, and metrics for evaluating similarity and clustering performance.

\section{Linear Algebra Distance Calculations}


\section{Similarity Metrics}

\subsection{Tanimoto Coefficient (also known as the Jaccard Index)}

\textbf{Range:} 0-1
\textbf{Higher is More Similar:} True

We define the Tanimoto Coefficient (TC) between two sets (A, B) as such:
\[TC(A,B) = TC(B,A) = \frac{|A\cap B|}{|A\cup B|}\ = \frac{|A\cap B|}{|A| + |B| - |A\cap B|}\]

The related Jaccard Distance is simply $1-TC$. 

This is a favorite for comparing bitmaps (i.e. chemical fingerprints), in part because calculating
the intersection and union of two bitvectors is very efficient.

As an example, given ligand A and ligand B with fingerprints:
\[A = |0|1|0|1|0|\]
\[B = |0|1|1|0|0|\]

\[TC(A,B) = \frac{|0|1|0|0|0|}{|0|1|1|1|0|} = \frac{1}{3} = 0.33 \]

\subsection{Dice-Sørensen Coefficient}

\textbf{Range:} 0-1
\textbf{Higher is More Similar:} True

Following the conventions above, the Dice-Sørensen Coefficient can be calculated by:
\[DC(A,B) = \frac{2|A\cap B|}{|A| + |B|}\]

Although similar to the Tanimoto Coefficient, this function does not satisfy the triangle inequality and is much less frequently used.

\section{Clustering Metrics}


\subsubsection{Silhouette Coefficient}
The silhouette coefficient is a measure of how similar an object is to its own cluster compared to other clusters.
First we calculate the mean intra-cluster distance $a_i$ by taking the average of each distance $d(i,j)$ for each object $i$ in cluster $C_i$:
\[
    a_i = \frac{1}{|C_I| - 1} \sum_{j \in C_I, i \neq j} d(i, j)
\]

Next we calculate the mean nearest-cluster distance $b_i$ for each object $i$ in cluster $C_i$:
\[
    b_i = \min_{J \neq I} \frac{1}{|C_J|} \sum_{j \in C_J} d(i, j)
\]

Then we can define the silhouette coefficient is defined for each object $i$ as:
\[
    SC_i = \frac{b_i - a_i}{\max(a_i, b_i)}
\]

Finally, by taking the mean silhouette coefficient over all objects, the final value will be between -1 and 1,
where a high value indicates a good clustering.

\end{document}